\section{Specialized frozen spin rings at JINR}
\label{sec:specialized-ring}

This section presents two specialized frozen spin lattices designed
exclusively for the search for the electric dipole moments of the proton and
the deuteron: an all-electrostatic ring for protons and a ring with combined
E+B deflectors for deuterons.

\subsection{Frozen spin option for the proton beam}

The frozen spin lattice for a proton beam is the simplest one in terms of
parameter choice and is uniquely defined. It consists of electrostatic
deflectors whose total length $L_{total}$ is determined by the beam energy and
the maximum permissible electrostatic field strength:
\begin{equation}
  L_{total} = 2\pi\,\frac{mc^2}{eE}\,\gamma\beta^2 .
  \label{eq:pring_length}
\end{equation}
For an energy of \SI{233}{MeV} and a maximum field of $E = \SI{120}{kV/cm}$,
the total length of all electrostatic deflectors is $L_{total} = \SI{219.55}{m}$.

With 50 deflectors placed around the ring, the length of a single deflector is
$L_{def} = \SI{4.39}{m}$. The 50 deflectors are divided into five groups, each
consisting of five FODO cells and containing ten deflectors
(Fig.~\ref{fig:pring_superperiod}).

\begin{figure}[htbp]
  \centering
  \includegraphics[width=0.8\textwidth]{pring_superperiod}
  \caption{Twiss functions of one superperiod of the electrostatic proton
    ring.}
  \label{fig:pring_superperiod}
\end{figure}

Figure~\ref{fig:pring_twiss} shows the whole ring, which consists of five
superperiods. Each superperiod comprises an achromat with drift spaces at both
ends, so that the phase advance per superperiod is not an integer multiple of
$2\pi$. The ring circumference is \SI{419.1}{m}. The layout of the
electrostatic ring is shown in Fig.~\ref{fig:pring_layout}; it closely
resembles a circle.

\begin{figure}[htbp]
  \centering
  \includegraphics[width=0.8\textwidth]{pring_twiss}
  \caption{Twiss functions along the whole electrostatic proton ring.}
  \label{fig:pring_twiss}
\end{figure}

\begin{figure}[htbp]
  \centering
  \includegraphics[width=0.6\textwidth]{pring_layout}
  \caption{Layout of the electrostatic proton ring.}
  \label{fig:pring_layout}
\end{figure}

Table~\ref{tab:pring} summarizes the ring parameters required to estimate its
dimensions and cost.

{\small
  \begin{longtable}{p{0.55\textwidth}p{0.38\textwidth}}
    \caption{Parameters of the frozen spin electrostatic ring for the proton
      EDM search.}
    \label{tab:pring} \\
    \toprule
    Parameter & Value \\
    \midrule
    \endfirsthead
    \toprule
    Parameter & Value \\
    \midrule
    \endhead
    Circumference, m & 419.1 \\
    Symmetry order (number of arcs and straight sections) & 5 \\
    Injection energy, MeV & 20 \\
    Proton magnetic anomaly $G$ & 1.793 \\
    Final energy, MeV & 233.8 \\
    EDM experiment energy, MeV & 233.8 \\
    Initial and final relative velocity and Lorentz factor
      & $\beta=0.2$, $\gamma=1.0213$; $\beta=0.6$, $\gamma=1.248$ \\
    Revolution frequency at the beginning and end of the cycle, MHz
      & 0.143--0.429 \\
    Harmonic number & 20 \\
    RF frequency at the beginning and end of the cycle, MHz & 2.8--8.6 \\
    \midrule
    \multicolumn{2}{c}{\textbf{Arcs}} \\
    \midrule
    Arc length, m & 70.1 \\
    Arc superperiodicity & regular lattice \\
    Arc lattice & 5 FODO cells \\
    Number of electrostatic deflectors per arc & 10 \\
    Deflector length, m & 4.39 \\
    Deflector type & cylindrical \\
    Maximum electric field of the deflector, MV/m & 12 \\
    Number of quadrupoles per arc & 9 + 1 (two end halves) \\
    Quadrupole type & electrostatic \\
    Quadrupole families & one focusing (LF), one defocusing (LD) \\
    Maximum quadrupole gradient in the arcs, kV/cm$^2$
      & LF $= 6.78$, LD $= -6.18$ \\
    Quadrupole length, m & 1.0 \\
    Number of chromatic sextupoles per arc & 10 \\
    Sextupole type & electric or magnetic \\
    Sextupole strengths & to be determined \\
    Betatron tunes $\nu_x/\nu_y$ & 6.12/4.72 (preliminary) \\
    Natural chromaticity $\xi_x/\xi_y$ & to be determined \\
    Transition energy, GeV & 3.1 \\
    Maximum $\beta_x/\beta_y$, m & 27.9/34.8 \\
    Maximum horizontal dispersion, m & 10.8 \\
    \midrule
    \multicolumn{2}{c}{\textbf{Straight sections between arcs}} \\
    \midrule
    Length of straight sections, m & $5\times 13.5$ \\
    Lattice of straight sections & DOFFOD \\
    Quadrupole gradients, T/m & QF $= 7.93$, QD $= -7.11$ \\
    Drift length between adjacent quadrupoles, m & 6.0 \\
    \bottomrule
  \end{longtable}
}

As a second option, a racetrack ring with two arcs and dispersion suppressors
at both ends of each arc was considered (Fig.~\ref{fig:pring_suppressor}). One
of the two options will be adopted in the course of further lattice
development, which will take spin decoherence into
account~\cite{Senichev:2013dra}.

\begin{figure}[htbp]
  \centering
  \includegraphics[width=0.9\textwidth]{pring_arc_suppressor}
  \caption{Arc with dispersion suppressors at both ends (racetrack option of
    the proton ring).}
  \label{fig:pring_suppressor}
\end{figure}

\subsection{Frozen spin option for the deuteron beam}

The frozen spin lattice for a deuteron beam differs from the proton one in the
composition of the equipment and in the elements used. Here the bending
deflector with E+B fields consists of electrostatic plates
generating a horizontal electric field, inserted into the gap of a dipole
magnet with a vertical magnetic field (Fig.~\ref{fig:dring_deflector}).

\begin{figure}[htbp]
  \centering
  \begin{panel}[b]{0.38\textwidth}
    \centering
    \includegraphics[width=\textwidth]{dring_deflector}
  \end{panel}
  \hfill
  \begin{panel}[b]{0.58\textwidth}
    \centering
    \includegraphics[width=\textwidth]{dring_magnet}
  \end{panel}
  \caption{Electrostatic deflector (a) and vertical-field dipole magnet with
    the electrostatic deflector inserted into its gap (b).}
  \label{fig:dring_deflector}
\end{figure}

The condition of stable motion on an orbit of radius $R$ follows from the
balance of the Lorentz and centrifugal forces:
\begin{equation}
  R = \frac{\gamma m_0 v^2}{e\left|\vec E + c\vec\beta\times\vec B\right|},
  \label{eq:dring_radius}
\end{equation}
from which the momentum precession frequency in the combined fields is
\begin{equation}
  \vec\Omega_p = \frac{e\vec E}{m_0\gamma\beta c} + \frac{e\,c\vec\beta\times\vec B}{m_0\gamma\beta c}.
\end{equation}
Subtracting $\vec\Omega_p$ from $\vec\Omega_{mdm}$ gives the spin precession
frequency relative to the momentum:
\begin{equation}
  \vec\Omega^{p}_{mdm} = \frac{e}{m_0\gamma}\left\{\gamma G\vec B
    + \left(\frac{1}{\gamma^2-1} - G\right)\frac{\vec\beta\times\vec E}{c}\right\}.
\end{equation}

A frozen spin lattice for EDM studies has to satisfy two conditions. The first
relates the external fields to the orbit radius, Eq.~\eqref{eq:dring_radius}
(see also~\cite{Farley:2004hp}). The second requires the spin precession
frequency relative to the direction of motion to vanish,
$\vec\Omega^{p}_{mdm} = 0$, i.e.,
$\gamma G\vec B = -\left(\frac{1}{\gamma^2-1}-G\right)\frac{\vec\beta\times\vec E}{c}$.
This fixes the relation between the fields,
\begin{equation}
  \vec E = -\frac{G\vec B c\gamma^2\beta^2}{\beta\left(1-\gamma^2\beta^2 G\right)},
\end{equation}
and the radius of curvature of the deflector,
\begin{equation}
  R = \frac{\gamma m_0 v^2}{eE\left(1-\dfrac{1}{G\gamma^2}\right)}.
\end{equation}
For an energy of \SI{270}{MeV}, the radius of curvature of the deflector is
$R = \SI{9.207}{m}$, and the total length of all deflectors is
$L_\Sigma = 2\pi R = \SI{57.85}{m}$. With 32 deflectors in total, the length of
a single deflector is $L_{def} = L_\Sigma/32 = \SI{1.808}{m}$.

Figure~\ref{fig:dring_vector} shows the vector diagram that defines the
relative directions of $\vec E$ and $\vec B$. The determining factor is the
positive vertical direction of the magnetic field, which provides the main
compensation of the centrifugal force. Since the deuteron magnetic anomaly is
negative, $G = -0.143$, the axis around which $\vec S$ precesses points
downwards. A zero precession frequency then requires the vector
$\vec\beta\times\vec E$ to point upwards, i.e., the electric field must be
directed as shown in Fig.~\ref{fig:dring_vector}.

\begin{figure}[htbp]
  \centering
  \includegraphics[width=0.75\textwidth]{dring_vector_diagram}
  \caption{Vector diagram of the fields acting on a polarized deuteron.}
  \label{fig:dring_vector}
\end{figure}

Based on these relations, a frozen spin ring lattice was developed.
Figure~\ref{fig:dring_arc} shows the lattice of one arc for the deuteron beam.
The arc consists of two successive achromats with zero dispersion at its middle
and at its ends and comprises eight FODO cells.

\begin{figure}[htbp]
  \centering
  \includegraphics[width=0.85\textwidth]{dring_arc_twiss}
  \caption{Twiss functions along an arc of the frozen spin deuteron ring.}
  \label{fig:dring_arc}
\end{figure}

As in the electrostatic proton ring, the dispersion modulation introduced in
the deuteron lattice provides suitable locations for sextupoles suppressing the
spin chromaticity. Figure~\ref{fig:dring_twiss} shows the Twiss functions of
the whole ring with a circumference of \SI{146}{m}; its layout is shown in
Fig.~\ref{fig:dring_layout}. The main parameters of the synchrotron for the
deuteron EDM search at \SI{270}{MeV} per nucleon are listed in
Table~\ref{tab:dring}.

\begin{figure}[htbp]
  \centering
  \includegraphics[width=0.85\textwidth]{dring_twiss}
  \caption{Twiss functions along the frozen spin deuteron ring.}
  \label{fig:dring_twiss}
\end{figure}

\begin{figure}[htbp]
  \centering
  \includegraphics[width=0.6\textwidth]{dring_layout}
  \caption{Layout and dimensions of the frozen spin deuteron ring.}
  \label{fig:dring_layout}
\end{figure}

{\small
  \begin{longtable}{p{0.55\textwidth}p{0.38\textwidth}}
    \caption{Parameters of the frozen spin ring with E+B
      deflectors for the deuteron EDM search.}
    \label{tab:dring} \\
    \toprule
    Parameter & Value \\
    \midrule
    \endfirsthead
    \toprule
    Parameter & Value \\
    \midrule
    \endhead
    Circumference, m & 145.85 \\
    Symmetry order (number of arcs and straight sections) & 2 \\
    Injection energy, MeV/u & 40 \\
    Final energy, MeV/u & 270.0 \\
    EDM experiment energy, MeV/u & 270.0 \\
    Deuteron magnetic anomaly $G$ & $-0.143$ \\
    Initial and final relative velocity and Lorentz factor
      & $\beta=0.2$, $\gamma=1.0213$; $\beta=0.485$, $\gamma=1.143$ \\
    Revolution frequency at the beginning and end of the cycle, MHz
      & 0.143--0.99 \\
    Harmonic number $h$ & 50 \\
    RF frequency at the beginning and end of the cycle, MHz & 7.15--49.9 \\
    Length of the quarter-wave RF cavity, m & 1.502 \\
    Momentum acceptance of the separatrix at maximum energy & $1.31\times10^{-3}$ \\
    Relative momentum spread in the bunch & $5\times10^{-4}$ \\
    \midrule
    \multicolumn{2}{c}{\textbf{Arcs}} \\
    \midrule
    Arc length, m & 51.27 \\
    Arc superperiodicity & 2 \\
    Arc lattice & 4 + 4 FODO cells \\
    Number of E+B deflectors per arc & 16 \\
    Deflector length, m & 1.80 \\
    Deflector type & combined magnetic and electric \\
    Maximum electric field of the deflector, MV/m & 12 \\
    Number of quadrupoles per arc & 15 + 1 (two end halves) \\
    Quadrupole type & magnetic multipole \\
    Quadrupole families & one focusing (QF), one defocusing (QD) \\
    Maximum quadrupole gradient in the arcs, kG/cm
      & QF $= 1.359$, QD $= -1.00$ \\
    Quadrupole length, m & 0.1 \\
    Number of chromatic sextupoles per arc & 16 \\
    Sextupole type & magnetic \\
    Sextupole strengths, kG/cm$^2$ & SF $= 0.12$, SD $= 0.15$ \\
    Betatron tunes $\nu_x/\nu_y$ & 4.70418/2.6140 (preliminary) \\
    Natural chromaticity $\xi_x/\xi_y$ & $-4.91629/{-3.99402}$ \\
    Transition energy, GeV & 5.3 \\
    Maximum $\beta_x/\beta_y$, m & 15.1/16.5 \\
    Maximum horizontal dispersion, m & 2.5 \\
    \midrule
    \multicolumn{2}{c}{\textbf{Straight sections between arcs}} \\
    \midrule
    Length of straight sections, m & $2\times 23.3$ \\
    Lattice of straight sections & $3\times$FODO \\
    Quadrupole gradients, kG/cm & QF $= 0.734$, QD $= -0.782$ \\
    Free drift spaces between adjacent quadrupoles, m & $12\times 3.2$ \\
    \bottomrule
  \end{longtable}
}

\subsection{Summary}

Two specialized frozen spin lattices have been developed for the search for
the proton and deuteron electric dipole moments. These results form the basis
for the physical design of the experimental facilities and their cost
estimate. In contrast to the quasi-frozen spin options of
Sections~\ref{sec:nica-bypass} and~\ref{sec:nuclotron}, these rings are
intended exclusively for experiments on the violation of fundamental
symmetries.

The material of this section has been reported
in~\cite{Palamarchuka:2026pyv,Kolokolchikov:2026xbv}.
